% cuRAMSES III — Avengers Assembly:
% non-standard cosmological physics suite
% (SIDM / dark energy & non-GR gravity / FDM / PBH)
%
% Title/subtitle per request:
%   제목: 어벤저스 어셈블리  부제: 뭘 좋아할지 몰라서 다 넣었어
%
\documentclass[fleqn,usenatbib]{mnras}

\usepackage{amsmath,amssymb}
\usepackage{graphicx}
\usepackage{booktabs}
\usepackage{multirow}
\usepackage{xcolor}
\usepackage{hyperref}

\newcommand{\ramses}{\textsc{ramses}}
\newcommand{\curamses}{\textsc{cuRAMSES}}
\newcommand{\camb}{\textsc{camb}}
\newcommand{\code}[1]{\texttt{#1}}
\newcommand{\Geff}{G_{\rm eff}}
\newcommand{\Mpl}{M_{\rm Pl}}
\newcommand{\kms}{\,{\rm km\,s^{-1}}}

\title[Avengers Assembly]
{Avengers Assembly: We Didn't Know What You'd Like,\\
So We Put Everything In}

\author[J.~Kim]{
Juhan Kim$^{1}$\thanks{E-mail: kjhan@kias.re.kr}
\\
$^{1}$Center for Advanced Computation, Korea Institute for Advanced Study,
85 Hoegiro, Dongdaemun-gu, Seoul 02455, Republic of Korea
}

\date{Accepted XXX. Received YYY; in original form ZZZ}
\pubyear{2026}

\begin{document}
\label{firstpage}
\pagerange{\pageref{firstpage}--\pageref{lastpage}}
\maketitle

% ===================================================================
\begin{abstract}
We present the non-standard cosmology suite of the \curamses\
adaptive mesh refinement (AMR) code (\curamses~III), which
assembles, within a single code base, essentially every popular
extension of $\Lambda$CDM: (i) \textbf{self-interacting dark
matter} (SIDM) with constant, velocity-dependent, anisotropic and
inelastic (two-state) cross-sections, implemented as cell-based
Monte Carlo scattering with machine-precision momentum and energy
conservation; (ii) a \textbf{dark energy and non-GR gravity}
battery comprising CPL dark energy with clustering perturbations,
massive-neutrino linear response, early dark energy, field-level
quintessence and purely kinetic k-essence, coupled quintessence,
parametrized quasi-static Horndeski gravity, Hu--Sawicki $f(R)$,
nDGP, symmetron, the environmentally damped dilaton, the
parameter-free tracker cubic Galileon, and QUMOND --- all built on
a unified Poisson-solver framework in which model corrections
compose multiplicatively, with an FFT-accelerated Newton stage for
the nonlinear scalar fields; (iii) \textbf{fuzzy dark matter}
(FDM), a hybrid Schr\"odinger--Poisson integrator with a spectral
kinetic operator on the base level and finite differences on
refined levels; and (iv) \textbf{evaporating primordial black
holes} (PBH), a table-driven global mass-function scheme with
optional local heating. Every dark-energy and gravity module has
been verified at the equation level against the literature, at the
solver level against analytic linear responses (to $10^{-5}$ or
better), and end-to-end with a controlled $64^3$ simulation suite
compared against an initial-conditions-consistent linear forward
model, agreeing to $2$--$5$ per cent with residual signs explained
by screening. We document the algorithms, unit conventions,
namelist interfaces and the verification methodology, providing a
self-contained reference for non-standard cosmological simulations
with AMR codes.
\end{abstract}

\begin{keywords}
methods: numerical -- cosmology: simulations -- dark energy --
dark matter -- gravitation -- neutrinos
\end{keywords}

% ===================================================================
\section{Introduction}
\label{sec:intro}

The $\Lambda$CDM concordance model provides an excellent fit to a
wide range of cosmological observations \citep{Planck2020}, yet the
physical nature of both dark components remains unknown. On the
dark-energy side, the cosmological constant may be replaced by a
dynamical field with an evolving equation of state
\citep{ChevallierPolarski2001,Linder2003}, by an early dark energy
component \citep{Poulin2019}, or the acceleration may signal a
breakdown of General Relativity on cosmological scales, as in
$f(R)$ gravity \citep{HuSawicki2007}, braneworld models
\citep{Schmidt2009}, Galileon \citep{Barreira2013}, symmetron and
dilaton theories \citep{HinterbichlerKhoury2010,Brax2010}, or MOND
\citep{Milgrom1983}. On the dark-matter side, small-scale tensions
motivate self-interacting dark matter
\citep{SpergelSteinhardt2000,TulinYu2018}, ultralight (fuzzy) dark
matter \citep{Hui2017,Schive2014}, and primordial black holes
\citep{Carr2021}, while neutrino oscillations guarantee a small
but non-zero hot component \citep{LesgourguesPastor2006}.

Stage-IV surveys will constrain all of these at the per-cent
level, and their nonlinear signatures can only be predicted by
simulations. Historically each extension has lived in its own
specialised code; comparing models then convolves physical
differences with code systematics. This paper describes the
opposite approach: one AMR code, \curamses\ \citep{Kim2026},
carrying \emph{all} of the above extensions behind uniform
namelist switches, with shared infrastructure (domain
decomposition, Morton-key octree, multigrid/CG/FFT Poisson
solvers, OpenMP+MPI parallelism) and a common verification
methodology. Hence the title\footnote{From the Korean
``\emph{mwol joahalji mollaseo da neoeosseo}'' --- we didn't
know what you'd like, so we put everything in.}.

The paper is organised in four physics sections, one per family:
Section~\ref{sec:sidm} describes self-interacting dark matter;
Section~\ref{sec:denongr} the dark energy and non-GR gravity
battery, including the unified Poisson framework, the nonlinear
scalar-field solvers and their verification; Section~\ref{sec:fdm}
fuzzy dark matter; and Section~\ref{sec:pbh} evaporating
primordial black holes. Section~\ref{sec:conclusion} concludes.
Throughout, code quantities use the supercomoving variables of
\citet{MartelShapiro1998}: $\tilde x = x_{\rm com}/(aB)$
(with $B$ the box size), $\tilde\rho = \rho a^3/(\rho_c\Omega_m)$
(so $\bar{\tilde\rho}=1$), $\tilde\phi = a^2\phi/(BH_0)^2$ and
${\rm d}\tilde t = H_0\,{\rm d}t/a^2$.

% ===================================================================
\section{Self-Interacting Dark Matter}
\label{sec:sidm}

\subsection{Physical motivation}

SIDM models introduce a non-zero scattering cross-section
$\sigma/m$ that allows dark matter particles to exchange momentum
and energy, potentially resolving the core--cusp, missing
satellites and too-big-to-fail problems
\citep{SpergelSteinhardt2000,TulinYu2018}. Astrophysical
constraints span dwarf galaxies ($v\sim30$--$100\kms$) to clusters
($v\sim10^3$--$3\times10^3\kms$); a constant
$\sigma/m\sim1\,{\rm cm^2\,g^{-1}}$ fits dwarf scales but may
violate cluster bounds, motivating velocity-dependent
cross-sections from light-mediator models
\citep{Feng2009,Tulin2013,Kaplinghat2016}.

\subsection{Monte Carlo scattering algorithm}

We implement cell-based Monte Carlo scattering following the
standard mesh/SPH approach \citep{Rocha2013,Robertson2017}. In
every AMR leaf cell, dark matter particles are collected from the
particle linked list, shuffled with the Fisher--Yates algorithm,
and paired sequentially. For a pair with relative velocity
$v_{\rm rel}$ the scattering probability per step is
\begin{equation}
P \;=\; \frac{\sigma/m(v)\,\bar m\,v_{\rm rel}\,\Delta t}
             {V_{\rm cell}}\,(N_{\rm dm}-1),
\label{eq:sidm_prob}
\end{equation}
with $\bar m$ the mean particle mass, $V_{\rm cell}$ the physical
cell volume and the factor $(N_{\rm dm}-1)$ accounting for all
possible partners. A scatter occurs when a uniform deviate falls
below $P$; post-scattering velocities are constructed in the
centre-of-mass frame,
\begin{equation}
\bmath{v}'_{i} = \bmath{v}_{\rm cm}
 + \frac{m_j}{m_i+m_j}\,|v'_{\rm rel}|\,\hat{\bmath{n}},\qquad
\bmath{v}'_{j} = \bmath{v}_{\rm cm}
 - \frac{m_i}{m_i+m_j}\,|v'_{\rm rel}|\,\hat{\bmath{n}},
\end{equation}
which conserves momentum by construction and kinetic energy
exactly for elastic scattering.

\subsection{Cross-sections and angular distributions}

Three velocity dependences are supported
(\code{sidm\_type}): constant $\sigma/m=\sigma_0$; Yukawa
(Born-regime) $\sigma(v)/m = \sigma_0[1+(v/v_0)^2]^{-2}$; and a
generic power law $\sigma(v)/m=\sigma_0(v/v_0)^n$. The angular
distribution is isotropic by default, or forward-peaked
Rutherford-like,
${\rm d}\sigma/{\rm d}\Omega \propto (1-\cos\theta+2\varepsilon)^{-2}$
\citep{Kahlhoefer2014}, sampled by CDF inversion.

\subsection{Inelastic (two-state) SIDM}

A two-state model \citep{TuckerSmithWeiner2001,Vogelsberger2019}
adds a mass splitting $\delta$: ground--ground pairs can
up-scatter ($g+g\to e+e$, allowed when
$\tfrac12\mu v_{\rm rel}^2>2\delta$), excited pairs down-scatter
exothermically, and mixed pairs scatter elastically; the relative
speed is adjusted to conserve total (kinetic + internal) energy.
An initial excited fraction is assigned at start-up in a
restart-safe manner. A dissipative channel
(\code{sidm\_fdiss}) and a DM--baryon drag term
(\code{sidm\_baryon}) are also available.

\subsection{Timestep, threading and diagnostics}

The maximum per-level probability $P_{\rm max}$ is fed back into
the timestep as
$\Delta t_{\rm new}\le (C_{\rm SIDM}/P_{\rm max})\,\Delta t_{\rm old}$
(default $C_{\rm SIDM}=0.1$). The sweep is OpenMP-parallel over
grids with per-thread decorrelated RNG streams and deterministic
reseeding, and cumulative momentum/energy residuals are reduced
across MPI ranks each step, providing runtime verification of the
conservation properties. All parameters live in
\code{\&SIDM\_PARAMS}.

% ===================================================================
\section{Dark Energy and non-GR Gravity}
\label{sec:denongr}

This section describes the largest part of the suite: eleven
dark-energy and modified-gravity models sharing one solver
infrastructure. Table~\ref{tab:models} summarises the models,
their parameters and their verification status.

\begin{table*}
\centering
\caption{The dark energy and gravity battery. ``Linear response''
lists the analytically validated unscreened limit; the last column
the accuracy achieved in dedicated unit tests of the discretised
equations (see Section~\ref{sec:verification}).}
\label{tab:models}
\begin{tabular}{lllll}
\toprule
Model & Switch / namelist & Background & Linear response & Unit-test accuracy \\
\midrule
CPL DE + perturbations & \code{w0,wa,cs2\_de,de\_table} & $w=w_0+w_a(1-a)$ & Sapone--Kunz closure & $\le0.1$--$0.8\%$ vs \camb\ ($z\ge0.5$) \\
Massive $\nu$ linear response & \code{use\_neutrino} & --- & $1+(\Omega_\nu/\Omega_{cb})R_\nu(k,a)$ & table-exact \\
Early dark energy & \code{use\_ede} & Poulin+19 fluid & smooth & analytic dilution exact \\
Quintessence & \code{use\_quintessence} & Klein--Gordon (shooting) & smooth ($c_s^2=1$) & $w$-consistency $2\times10^{-5}$ \\
k-essence & \code{use\_kessence} & $P(X)=-X+X^2$ (exact) & Helmholtz with $c_s^2(a)$ & analytic limits exact \\
Coupled quintessence & \code{use\_coupled\_de} & KG + mass evolution & $\Geff=G(1+2\beta^2)$ & isolation runs $\le1\%$ \\
Horndeski $\mu(a,k)$ & \code{use\_horndeski} & $\Lambda$CDM/CPL & $\mu=1+\mu_0\Omega_{\rm DE}(a)/\Omega_{\rm DE,0}$ & exact by construction \\
$f(R)$ Hu--Sawicki & \code{use\_fR} & $\Lambda$CDM & $\mu=1+\frac13\frac{k^2}{k^2+m^2a^2}$ & $6\times10^{-7}$ \\
nDGP & \code{use\_nDGP} & $\Lambda$CDM & $\Geff/G=1+1/(3\beta)$ & force field $1.5\times10^{-6}$ \\
Symmetron & \code{use\_symmetron} & $\Lambda$CDM & $F_5/F_N=2\beta^2\bar\chi^2$ & exact (Brax+12 eq.~68 match) \\
Dilaton (Brax+12) & \code{use\_dilaton} & $\Lambda$CDM & $2\beta(a)^2 k^2/(k^2+m^2a^2)$ & 5 digits \\
Cubic Galileon (tracker) & \code{use\_galileon} & tracker $E^2(a)$ & $\Geff/G(a{=}1)=1.844$ & 5 digits \\
QUMOND & \code{use\_mond} & --- & $\nu(|g_N|/a_0)$ & $\mu\nu$ conjugacy $2\times10^{-16}$ \\
\bottomrule
\end{tabular}
\end{table*}

\subsection{Unified Poisson-solver framework}
\label{sec:unified}

In supercomoving units the code solves
$\tilde\nabla^2\tilde\phi = \tfrac32\Omega_m a(\tilde\rho-1)$.
All smooth-dark-sector corrections enter this equation as
multiplicative factors on the source. A single function,
\code{cosmo\_poisson\_fourpi()}, provides the prefactor
consistently to the multigrid, conjugate-gradient and coarse-level
solvers, and the FFT solvers apply the same corrections per
Fourier mode in the Green's function, so that scale-dependent
responses cost nothing extra. The factors compose:
\begin{equation}
\frac32\Omega_{cb}\,a
\;\times\;
\underbrace{\bigl[1+\tfrac{\Omega_\nu}{\Omega_{cb}}R_\nu(k,a)\bigr]}_{\rm neutrinos}
\underbrace{\bigl[1+\tfrac{\Omega_{\rm DE}(a)}{\Omega_{cb}}R_{\rm DE}(k,a)\bigr]}_{\rm DE\ clustering}
\underbrace{\vphantom{\bigl[\tfrac00\bigr]}\mu(a,k)}_{\rm Horndeski}
\underbrace{\vphantom{\bigl[\tfrac00\bigr]}e^{-\beta(\varphi-\varphi_0)}}_{\rm coupled\ DE}.
\end{equation}

\subsection{CPL dark energy with clustering perturbations}
\label{sec:cpl}

The background uses the exact CPL integral
$f_{\rm DE}(a)=a^{-3(1+w_0+w_a)}e^{-3w_a(1-a)}$ in the
supercomoving Friedmann factors, the growth factor and the initial
velocity scaling. Perturbations are supported in two modes:

\emph{(i) Table mode.} A \camb-derived table of the effective
ratio $R_{\rm DE}(k,a)=\delta_{\rm DE}/\delta_m$ is measured with
the Weyl-potential method (the ratio of potential per unit matter
contrast between a clustering-DE and a smooth-DE universe), which
yields the Newtonian-gauge response appropriate for the Poisson
source; synchronous-gauge $\delta_{\rm DE}$ would over-correct by
an order of magnitude.

\emph{(ii) Quasi-static fallback.} When no table is supplied and
$c_s^2>0$, the code applies the \citet{SaponeKunz2009} closure
\begin{equation}
\frac{\delta_{\rm DE}}{\delta_m}
= \frac{1+w}{(1-3w) + \dfrac{2c_s^2 k^2}{3a^2H^2}},
\qquad
\mbox{i.e. } {\rm d.f.} = 1+\frac{\alpha}{k^2+\kappa^2},
\label{eq:sk}
\end{equation}
with $\kappa^2=\tfrac32(1-3w)a^2E^2(L/\lambda_H)^2/c_s^2$ and
$\alpha=(\Omega_\Lambda f_{\rm DE}a^3/\Omega_m)(1+w)\,
\tfrac32 a^2E^2 (L/\lambda_H)^2/c_s^2$. Correct limits: a
$k$-independent boost $1+(\rho_{\rm DE}/\rho_m)(1+w)/(1-3w)$ for
$c_s^2\to0$, unity below the sound horizon, and bit-wise unity for
$\Lambda$. Validated against \camb\ (same Weyl method) to
$\le0.1$ per cent for $(w_0,w_a)=(-0.9,0.1)$ and $\le0.8$ per cent
for $w_0=-0.8$ at $z\ge0.5$ ($\sim4$ per cent at $z=0$, where the
quasi-static approximation degrades --- use the table mode for
precision at low redshift). An earlier release used a closure with
the opposite sign (a $15$--$29$ per cent source suppression);
results obtained with it should be discarded.

\subsection{Massive neutrino linear response}

Neutrino perturbations are added without neutrino particles:
$\delta_\nu(k,a) = R_\nu(k,a)\,\delta_{cb}(k,a)$ with
$R_\nu$ precomputed from \camb\ \citep{Lewis2000}, entering the
Poisson source as in Section~\ref{sec:unified}. This avoids
neutrino shot noise entirely and is exact at linear order
\citep[cf.][]{AliHaimoudBird2013}.

\subsection{Early dark energy}

The \citet{Poulin2019} fluid form
$\rho_{\rm EDE}(a)=\omega_{\rm ede}\,\rho_{\rm std}(a_c)\,
2/[1+(a/a_c)^{3(1+w_{\rm ede})}]$ is added to the Friedmann
factors. The implementation reproduces the analytic dilution slope
$-3(1+w)$ and the analytic peak
$f_{\rm EDE}^{\rm max}=1.14\,\omega_{\rm ede}$ at $a=3^{1/4}a_c$
exactly. Two conventions matter when mapping CMB constraints:
$\omega_{\rm ede}$ is the fraction at $a_c$ (not the peak), and
the background carries no radiation, so at $z_c\sim3000$ the
fraction differs from Planck-style $f_{\rm EDE}$ by $\sim2\times$.

\subsection{Scalar-field dark energy: quintessence and k-essence}
\label{sec:sde}

A dedicated module tabulates scalar-field backgrounds on a
$\ln a$ grid ($a=10^{-6}$--$1.55$) with lazy, self-healing
initialisation (the tables rebuild automatically if the cosmology
is overridden by IC headers).

\emph{Quintessence} \citep{RatraPeebles1988}: the Klein--Gordon
and Friedmann equations are integrated with RK4 for
Ratra--Peebles ($V\propto\varphi^{-\alpha}$) or exponential
potentials, the amplitude fixed by shooting to
$\Omega_\varphi(a{=}1)=\Omega_\Lambda$. The tabulated $w(a)$ obeys
${\rm d}\ln f_{\rm DE}/{\rm d}\ln a=-3(1+w)$ to $2\times10^{-5}$,
and the RP tracker attractor $w=-2/(\alpha+2)$ is recovered in the
matter era.

\emph{k-essence} \citep{ArmendarizPicon2001}: for the purely
kinetic $P(X)=-X+X^2$ the shift-symmetry integral
$P_{,X}\sqrt{2X}\,a^3={\rm const}$ gives $X(a)$ algebraically;
$w\to1/3$ and $c_s^2\to1/3$ at early times (a dark-radiation-like
component) and $w\to-1$ today. The model's $c_s^2(a)$ feeds the
Helmholtz closure of equation~(\ref{eq:sk}) automatically.

\subsection{Coupled quintessence}
\label{sec:cde}

Conformally coupled dark matter, $m_{\rm dm}\propto
e^{-\beta\varphi/\Mpl}$ \citep{Amendola2000}, adds four effects,
all implemented and individually verified: (i) the modified
expansion, including the DM mass-evolution factor in the
\emph{matter} term of the Friedmann integrals (its omission is
detectable: an isolation run reproduced the inconsistent-$H$
growth prediction to $0.2$ per cent); (ii) the same factor
multiplying the Poisson source; (iii) the momentum-conservation
velocity term $+\beta\dot\varphi\,\bmath{v}$, split across the two
half-kicks of the KDK integrator; and (iv) the fifth-force
enhancement $\Geff=G(1+2\beta^2)$ between matter particles.

\subsection{Parametrized Horndeski gravity}

The quasi-static, scale-independent limit of Horndeski gravity is
implemented as
$\mu(a)=1+\mu_0\,\Omega_{\rm DE}(a)/\Omega_{\rm DE,0}$
multiplying the Poisson source, with an optional Compton-mass
$k$-dependence
$\mu(a,k)=1+[\mu(a)-1]\,k^2/(k^2+a^2m^2)$ applied exactly in the
FFT solver paths (the multigrid path uses the $k\to\infty$ limit).

\subsection{Screened modified gravity: the scalar-field solvers}
\label{sec:mg}

Five models solve a nonlinear elliptic equation for an extra
scalar degree of freedom on the AMR hierarchy, with the fifth
force added to the particle acceleration.

\subsubsection{$f(R)$ Hu--Sawicki}
The quasi-static equation in code units is
\begin{equation}
\tilde\nabla^2 f_R
= \frac{a^2}{3}\Bigl(\frac{H_0 B}{c}\Bigr)^{\!2}
  \bigl[\tilde R(f_R)-\bar{\tilde R}\bigr]
 -\Omega_m \Bigl(\frac{H_0 B}{c}\Bigr)^{\!2}\frac{\tilde\rho-1}{a},
\end{equation}
with the Hu--Sawicki inversion
$\tilde R = \bar{\tilde R}_0 (f_{R0}/f_R)^{1/(n+1)}$ and force
$F_5=+\tfrac{a^2}{2}(c/H_0B)^2\tilde\nabla f_R$
\citep{HuSawicki2007,Oyaizu2008,Li2012}. The linearized mass term
is a Yukawa (the Jacobian is negative definite) and the chameleon
limit $\tilde R\to\bar{\tilde R}+3\Omega_m\delta/a^3$ is exact.
The discretised linear response reproduces
$\mu(k)-1=\tfrac13 k^2/(k^2+m^2a^2)$ to $6\times10^{-7}$ across
$a\in[0.25,1]$, $f_{R0}\in\{-10^{-5},-10^{-6}\}$, $n\in\{1,2\}$.

\subsubsection{nDGP}
The brane-bending mode obeys the Vainshtein equation
\citep{Schmidt2009,Li2013}; in code units
\begin{equation}
\tilde\nabla^2\varphi
+\frac{1}{12\,\Omega_{rc}\beta(a)\,a^4}
\bigl[(\tilde\nabla^2\varphi)^2-(\tilde\nabla_i\tilde\nabla_j\varphi)^2\bigr]
= \frac{\Omega_m a}{\beta(a)}(\tilde\rho-1),
\end{equation}
with $\beta=1+2Hr_c(1+\dot H/3H^2)$ and $F_5=-\tfrac12\tilde\nabla\varphi$,
so that the unscreened limit is $\Geff/G=1+1/(3\beta)$
\citep{Winther2015}.

\subsubsection{Symmetron}
The dimensionless field $\chi=\varphi/\varphi_\star$ obeys
$\tilde\nabla^2\chi = \tfrac{a^2}{2\tilde\lambda^2}
[\tilde\rho\,(a_{\rm ssb}/a)^3-1]\chi
+\tfrac{a^2}{2\tilde\lambda^2}\chi^3$
\citep{HinterbichlerKhoury2010,Davis2012}, identical to eq.~(68)
of \citet{Brax2012}, with force
$F_5=-6\Omega_m\beta^2(\lambda/B)^2(a^2/a_{\rm ssb}^3)\,
\chi\tilde\nabla\chi$, whose unscreened limit is the canonical
$F_5/F_N=2\beta^2\bar\chi^2$. Because $\chi=0$ is an exact fixed
point of the homogeneous equation, the broken-phase VEV
$\bar\chi(a)=\sqrt{1-(a_{\rm ssb}/a)^3}$ is seeded into
uninitialised cells at every solve (this also covers restarts and
newly refined cells).

\subsubsection{Environmentally damped dilaton}
The original $r=3/2$ dilaton of \citet{Brax2010}, with
$A(\varphi)=1+\tfrac{A_2}{2}\varphi^2/\Mpl^2$,
$m^2(a)=3A_2H^2(a)$ and $\beta(a)=\beta_0 a^{3\Omega_m}$
\citep{Brax2012}, is solved in $\chi=\varphi/\Mpl$:
\begin{equation}
\tilde\nabla^2\chi
= \frac{3\Omega_m A_2 B_2}{a}\,(\tilde\rho\,\chi-\bar\chi)
+ a^2 B_2\bigl[v_\varphi(\chi)-v_\varphi(\bar\chi)\bigr],
\end{equation}
with $v_\varphi(\chi)=-3\Omega_m\beta_0(A_2\chi/\beta_0)^{1-3/s}$,
$s=3\Omega_m$, $B_2=(H_0B/c)^2$ and
$F_5=-(c/H_0B)^2a^2A_2\,\chi\tilde\nabla\chi$. The exponent
$1-3/s<0$ makes the Newton Jacobian negative definite. The
background is an exact fixed point of the discrete operator, the
linear response matches $2\beta(a)^2k^2/(k^2+m^2a^2)$ to five
digits, and a $\delta=10^5$ top-hat drives $\chi\to0.03\bar\chi$
(Damour--Polyakov screening). User parameters are simply
$\beta_0$ and the present force range
$\lambda = 2998\,\xi\,h^{-1}$Mpc (then $A_2=1/3\xi^2$).

\subsubsection{Tracker cubic Galileon}
On the tracker, $H\dot\varphi=\xi H_0^2\Mpl$ with
$c_2=-6c_3\xi$, the model is parameter-free:
$\xi=\sqrt{6(1-\Omega_m)}$,
$E^2(a)=\tfrac12[\Omega_ma^{-3}+\sqrt{\Omega_m^2a^{-6}+4(1-\Omega_m)}]$
(wired through the background dispatch, so the modified expansion
and growth are automatic; $w_0\simeq-1.18$), and the perturbation
coefficients of \citet{Barreira2013} reduce to
$\beta_1=\tfrac{\xi}{3}[2\dot H/H^2-1+(1-\Omega_m)/E^4]$,
$\beta_2=2E^2\beta_1/\xi^2$. The field equation is the nDGP
template with coefficient $1/(3\beta_1a^4)$ and source
$\Omega_m a\,\delta/\beta_2$; the Poisson back-reaction term
integrates to a potential proportional to the field, so the entire
fifth force is one gradient, $F_5=+\tfrac{\xi}{6E^2}\tilde\nabla u$.
The unscreened response matches the analytic
$-\xi/(9\beta_2E^2)$ to five digits at $a=0.25$--$1$:
$\Geff/G(a{=}1)=1.844$, decaying as $1/E^4$ into the past. Since
$\Geff/G-1<10^{-3}$ for $z\gtrsim2.5$, the solver is skipped at
early times.

\subsubsection{QUMOND}
The quasi-linear formulation of MOND \citep{Milgrom2010} defines a
curl-free field via
$\nabla^2\Phi=\nabla\cdot[\nu(|\nabla\Phi_N|/a_0)\nabla\Phi_N]$.
Both an algebraic force-correction mode and the full
phantom-density mode (compute $\rho_{\rm ph}=\nabla\cdot[(\nu-1)
\nabla\Phi_N]/4\pi G$, re-solve Poisson) are implemented, with the
simple and standard interpolation pairs (verified mutually
conjugate to $2\times10^{-16}$), an external-field option, and
$a_0$ converted to code units through the exact supercomoving
acceleration factor (so the MOND scale is correctly constant in
proper units at all redshifts). An AQUAL fixed-point mode is also
available \citep[cf.][]{Lughausen2015,Candlish2015}.

\subsection{Numerical solution of the scalar equations}
\label{sec:solver}

All five scalar solvers share one machinery: nonlinear
Newton--Gauss--Seidel relaxation with true three-dimensional
red--black colouring (cell parity $i{+}j{+}k$), Neumann closure of
the stencil at coarse--fine boundaries, one-sided force gradients
at level edges, and MPI ghost exchange each sweep. On the fully
refined domain level the long-wavelength modes --- which plain
relaxation cannot converge in $\mathcal{O}(10)$ sweeps --- are
solved spectrally: a Newton--Helmholtz stage
$(\tilde\nabla^2-\bar m^2)\,\delta u=-\mathcal{R}$ with the
level-averaged Jacobian mass for $f(R)$/symmetron/dilaton, and a
\citet{Winther2015}-style operator splitting for the
Vainshtein-type models (solve the cell-local quadratic for
$\tilde\nabla^2\varphi$, then invert the Laplacian by FFT). In a
controlled test the operator-split stage reached a lower residual
in five iterations than $2\times10^4$ damped relaxation sweeps.
Two MPI subtleties deserve mention as a warning to implementers:
every rank must enter the per-level ALLREDUCE of the convergence
test and the final force ghost-exchange \emph{even when it owns no
grids on that level}; early returns on empty ranks deadlock or
abort the run at the first sparsely refined level.

\subsection{Verification}
\label{sec:verification}

Verification proceeded on three levels.

\emph{Equation level.} Every model was re-derived from its
original reference (unit conversions included) and the code
compared term by term; this stage found and fixed, among others, a
wrong-sign DE-perturbation closure, a tachyonic (non-convergent)
$f(R)$ operator, a double-counted nDGP coupling, a symmetron
force normalisation missing the box-unit conversion, and a
zero-pinned symmetron field.

\emph{Solver level.} The discretised operators were exercised in
isolation against analytic linear responses; the accuracies are
listed in Table~\ref{tab:models}.

\emph{End-to-end level.} A controlled suite of $64^3$, $100\,
h^{-1}$Mpc dark-matter-only simulations was run from shared
Zel'dovich initial conditions ($z=49$, \camb\ spectrum) for
$\Lambda$CDM, quintessence, k-essence, Horndeski, coupled DE (plus
per-effect isolation runs), the tracker Galileon and the dilaton.
The measured $P_{\rm model}/P_{\Lambda{\rm CDM}}$ in the band
$k=0.05$--$0.12\,h\,{\rm Mpc}^{-1}$ at $z\simeq0$ was compared
against an \emph{initial-conditions-consistent} linear forward
model, which (i) rescales the initial displacements by each
model's own velocity conversion factor --- \ramses\ converts the
shared velocity files with the model's $v_{\rm fact}$, an effect
reaching $(1.077)^2$ for k-essence and $(1.109)^2$ for coupled DE
--- and (ii) propagates the exact $(\delta,\theta)$ initial vector
through each model's linear system. Results:

\begin{table}
\centering
\caption{End-to-end suite: simulated vs.\ predicted
$P/P_{\Lambda\rm CDM}$ ($k=0.05$--$0.12\,h/$Mpc, $z\simeq0$).}
\label{tab:sims}
\begin{tabular}{lccc}
\toprule
Model & Simulation & Forward model & Ratio \\
\midrule
Quintessence (RP $\alpha{=}1$) & 0.832 & 0.817 & 1.018 \\
k-essence ($X_0/M^4{=}0.5005$) & 0.875 & 0.881 & 0.993 \\
Horndeski ($\mu_0{=}0.2$)      & 1.036 & 1.052 & 0.985 \\
Coupled DE ($\beta{=}0.1$)     & 0.966 & 1.005 & 0.961 \\
Galileon (tracker)             & 1.148 & 1.208 & 0.950 \\
Dilaton ($\beta_0{=}0.5$, $\lambda{=}20$) & 1.039 & 1.071 & 0.970 \\
\bottomrule
\end{tabular}
\end{table}

All models agree to $2$--$5$ per cent; the Galileon and dilaton
deficits carry the expected sign, since Vainshtein and
Damour--Polyakov screening remove part of the fifth-force boost
already at quasi-linear wavenumbers \citep[cf.][]{Barreira2013}.
Each coupled-DE ingredient (mass evolution, friction, $\Geff$) was
additionally isolated with toggle runs and verified at the
$\sim1$ per cent level. The forward model, the IC generator and
all unit-test drivers ship with the code
(\code{patch/cuRamses/aux/}).

% ===================================================================
\section{Fuzzy Dark Matter}
\label{sec:fdm}

Ultralight scalar dark matter with
$m_a\sim10^{-22}\,$eV \citep{Hui2017} is evolved as a classical
field $\psi$ obeying the Schr\"odinger--Poisson system. The
integrator uses drift--kick--drift operator splitting per AMR
level: the kinetic half-steps apply
$\exp(-{\rm i}\hbar k^2\Delta t/4a^2 m_a)$ spectrally on the base
level (FFT) and an explicit finite-difference Laplacian on refined
levels, while the potential kick applies
$\exp(-{\rm i}m_a\Phi\Delta t/\hbar)$ in real space; a
fourth-order Yoshida composition is optional
(\code{fdm\_split\_order=4}). Refinement follows the de~Broglie
wavelength (a minimum of \code{fdm\_nrefine\_dB} cells per
$\lambda_{\rm dB}$, with a density floor), reproducing the
solitonic cores and interference fringes characteristic of FDM
haloes \citep{Schive2014}. A hybrid mode
(\code{fdm\_use\_hjm}) evolves a Hamilton--Jacobi--Madelung fluid
on coarse levels and switches to the wave solver above
\code{fdm\_first\_wave\_level}, and \code{fdm\_hybrid} allows
mixed FDM + particle (stars/sinks) runs. The wavefunction is
stored as $(\psi_{\rm re},\psi_{\rm im})$ on the AMR mesh with
$|\psi|^2$ the total DM density, and dedicated light-cone and
output modules are provided. A detailed description and validation
of the FDM solver is presented in a companion paper.

% ===================================================================
\section{Evaporating Primordial Black Holes}
\label{sec:pbh}

The PBH module models dark matter that is partly composed of
Hawking-evaporating primordial black holes \citep{Hawking1975,
Carr2021}. Rather than tracking per-particle masses, it uses one
global mass function: all DM particles share the mixed-mass factor
\begin{equation}
w(a) \;=\; 1-f_{\rm PBH} + f_{\rm PBH}\,g(a),
\end{equation}
where $g(a)$ is the surviving PBH mass fraction from a
precomputed evaporation table and $f_{\rm PBH}$
(\code{pbh\_fraction}) the initial PBH share of the dark matter.
Particle masses are rescaled exactly by table arithmetic --- there
is no per-particle state and no evaporation CFL constraint. The
energy released between steps follows the cumulative heating
kernel $\tilde Q(a)=\int q\,g\,{\rm d}t$; it is either deposited
as local heat in the gas (\code{pbh\_energy\_sink='local\_heat'},
optionally amplified by a linear-response boost factor) or removed
from the simulation (\code{'removed'}, appropriate for free-streaming
decay products). Restart provenance (table checksum,
normalisation epoch) is persisted so that a run cannot silently
continue with a different evaporation model. The table generator
supports arbitrary monochromatic or extended PBH mass functions.

% ===================================================================
\section{Conclusions}
\label{sec:conclusion}

We have assembled, in one AMR code, the four families of
non-standard cosmological physics most often invoked beyond
$\Lambda$CDM --- self-interacting dark matter, a comprehensive
dark-energy and modified-gravity battery, fuzzy dark matter, and
evaporating primordial black holes --- behind uniform namelist
switches and on shared solver infrastructure. Beyond the
implementations themselves, we highlight the verification
methodology: every model is checked at the equation level against
its original reference, at the solver level against analytic
linear responses, and end-to-end against an
initial-conditions-consistent forward model, with per-effect
isolation runs where a model bundles several couplings. We commend
this three-level protocol to implementers of non-standard physics
in any simulation code: in our experience each level catches
defects invisible to the others, ranging from wrong-sign
perturbation closures to factor-of-two integrator bookkeeping and
MPI collective mismatches that only manifest on sparsely refined
levels.

All modules recover $\Lambda$CDM bit-wise when disabled. The
namelist interfaces, verification scripts and the simulation-suite
tooling are distributed with the code.

\section*{Acknowledgements}
The author thanks the KIAS Center for Advanced Computation for
computing resources. Portions of the implementation and the
verification campaign were carried out with the assistance of
Claude (Anthropic).

\section*{Data Availability}
The code, namelist examples, verification drivers and the
$64^3$ test-suite configuration are available from the
\curamses\ repository. Simulation outputs will be shared on
reasonable request.

% ===================================================================
\begin{thebibliography}{99}

\bibitem[\protect\citeauthoryear{Ali-Ha\"imoud \& Bird}{2013}]{AliHaimoudBird2013}
Ali-Ha\"imoud Y., Bird S., 2013, MNRAS, 428, 3375

\bibitem[\protect\citeauthoryear{Amendola}{2000}]{Amendola2000}
Amendola L., 2000, Phys. Rev. D, 62, 043511

\bibitem[\protect\citeauthoryear{Armendariz-Picon et al.}{2001}]{ArmendarizPicon2001}
Armendariz-Picon C., Mukhanov V., Steinhardt P.~J., 2001, Phys. Rev. D, 63, 103510

\bibitem[\protect\citeauthoryear{Barreira et al.}{2013}]{Barreira2013}
Barreira A., Li B., Hellwing W.~A., Baugh C.~M., Pascoli S., 2013, J. Cosmology Astropart. Phys., 10, 027

\bibitem[\protect\citeauthoryear{Brax et al.}{2010}]{Brax2010}
Brax P., van de Bruck C., Davis A.-C., Shaw D., 2010, Phys. Rev. D, 82, 063519

\bibitem[\protect\citeauthoryear{Brax et al.}{2012}]{Brax2012}
Brax P., Davis A.-C., Li B., Winther H.~A., Zhao G.-B., 2012, J. Cosmology Astropart. Phys., 10, 002

\bibitem[\protect\citeauthoryear{Candlish et al.}{2015}]{Candlish2015}
Candlish G.~N., Smith R., Fellhauer M., 2015, MNRAS, 446, 1060

\bibitem[\protect\citeauthoryear{Carr \& K\"uhnel}{2021}]{Carr2021}
Carr B., K\"uhnel F., 2021, arXiv:2110.02821

\bibitem[\protect\citeauthoryear{Chevallier \& Polarski}{2001}]{ChevallierPolarski2001}
Chevallier M., Polarski D., 2001, Int. J. Mod. Phys. D, 10, 213

\bibitem[\protect\citeauthoryear{Davis et al.}{2012}]{Davis2012}
Davis A.-C., Li B., Mota D.~F., Winther H.~A., 2012, ApJ, 748, 61

\bibitem[\protect\citeauthoryear{Feng et al.}{2009}]{Feng2009}
Feng J.~L., Kaplinghat M., Tu H., Yu H.-B., 2009, J. Cosmology Astropart. Phys., 07, 004

\bibitem[\protect\citeauthoryear{Hawking}{1975}]{Hawking1975}
Hawking S.~W., 1975, Comm. Math. Phys., 43, 199

\bibitem[\protect\citeauthoryear{Hinterbichler \& Khoury}{2010}]{HinterbichlerKhoury2010}
Hinterbichler K., Khoury J., 2010, Phys. Rev. Lett., 104, 231301

\bibitem[\protect\citeauthoryear{Hu \& Sawicki}{2007}]{HuSawicki2007}
Hu W., Sawicki I., 2007, Phys. Rev. D, 76, 064004

\bibitem[\protect\citeauthoryear{Hui et al.}{2017}]{Hui2017}
Hui L., Ostriker J.~P., Tremaine S., Witten E., 2017, Phys. Rev. D, 95, 043541

\bibitem[\protect\citeauthoryear{Kahlhoefer et al.}{2014}]{Kahlhoefer2014}
Kahlhoefer F., Schmidt-Hoberg K., Frandsen M.~T., Sarkar S., 2014, MNRAS, 437, 2865

\bibitem[\protect\citeauthoryear{Kaplinghat et al.}{2016}]{Kaplinghat2016}
Kaplinghat M., Tulin S., Yu H.-B., 2016, Phys. Rev. Lett., 116, 041302

\bibitem[\protect\citeauthoryear{Kim}{2026}]{Kim2026}
Kim J., 2026, MNRAS, submitted (cuRAMSES I)

\bibitem[\protect\citeauthoryear{Lesgourgues \& Pastor}{2006}]{LesgourguesPastor2006}
Lesgourgues J., Pastor S., 2006, Phys. Rep., 429, 307

\bibitem[\protect\citeauthoryear{Lewis et al.}{2000}]{Lewis2000}
Lewis A., Challinor A., Lasenby A., 2000, ApJ, 538, 473

\bibitem[\protect\citeauthoryear{Li et al.}{2012}]{Li2012}
Li B., Zhao G.-B., Teyssier R., Koyama K., 2012, J. Cosmology Astropart. Phys., 01, 051

\bibitem[\protect\citeauthoryear{Li et al.}{2013}]{Li2013}
Li B., Zhao G.-B., Koyama K., 2013, J. Cosmology Astropart. Phys., 05, 023

\bibitem[\protect\citeauthoryear{Linder}{2003}]{Linder2003}
Linder E.~V., 2003, Phys. Rev. Lett., 90, 091301

\bibitem[\protect\citeauthoryear{Llinares et al.}{2014}]{Llinares2014}
Llinares C., Mota D.~F., Winther H.~A., 2014, A\&A, 562, A78

\bibitem[\protect\citeauthoryear{L\"ughausen et al.}{2015}]{Lughausen2015}
L\"ughausen F., Famaey B., Kroupa P., 2015, Canadian J. Phys., 93, 232

\bibitem[\protect\citeauthoryear{Martel \& Shapiro}{1998}]{MartelShapiro1998}
Martel H., Shapiro P.~R., 1998, MNRAS, 297, 467

\bibitem[\protect\citeauthoryear{Milgrom}{1983}]{Milgrom1983}
Milgrom M., 1983, ApJ, 270, 365

\bibitem[\protect\citeauthoryear{Milgrom}{2010}]{Milgrom2010}
Milgrom M., 2010, MNRAS, 403, 886

\bibitem[\protect\citeauthoryear{Oyaizu}{2008}]{Oyaizu2008}
Oyaizu H., 2008, Phys. Rev. D, 78, 123523

\bibitem[\protect\citeauthoryear{Planck Collaboration}{2020}]{Planck2020}
Planck Collaboration, 2020, A\&A, 641, A6

\bibitem[\protect\citeauthoryear{Poulin et al.}{2019}]{Poulin2019}
Poulin V., Smith T.~L., Karwal T., Kamionkowski M., 2019, Phys. Rev. Lett., 122, 221301

\bibitem[\protect\citeauthoryear{Ratra \& Peebles}{1988}]{RatraPeebles1988}
Ratra B., Peebles P.~J.~E., 1988, Phys. Rev. D, 37, 3406

\bibitem[\protect\citeauthoryear{Robertson et al.}{2017}]{Robertson2017}
Robertson A., Massey R., Eke V., 2017, MNRAS, 465, 569

\bibitem[\protect\citeauthoryear{Rocha et al.}{2013}]{Rocha2013}
Rocha M., Peter A.~H.~G., Bullock J.~S., et al., 2013, MNRAS, 430, 81

\bibitem[\protect\citeauthoryear{Sapone \& Kunz}{2009}]{SaponeKunz2009}
Sapone D., Kunz M., 2009, Phys. Rev. D, 80, 083519

\bibitem[\protect\citeauthoryear{Schive et al.}{2014}]{Schive2014}
Schive H.-Y., Chiueh T., Broadhurst T., 2014, Nature Physics, 10, 496

\bibitem[\protect\citeauthoryear{Schmidt}{2009}]{Schmidt2009}
Schmidt F., 2009, Phys. Rev. D, 80, 043001

\bibitem[\protect\citeauthoryear{Spergel \& Steinhardt}{2000}]{SpergelSteinhardt2000}
Spergel D.~N., Steinhardt P.~J., 2000, Phys. Rev. Lett., 84, 3760

\bibitem[\protect\citeauthoryear{Teyssier}{2002}]{Teyssier2002}
Teyssier R., 2002, A\&A, 385, 337

\bibitem[\protect\citeauthoryear{Tucker-Smith \& Weiner}{2001}]{TuckerSmithWeiner2001}
Tucker-Smith D., Weiner N., 2001, Phys. Rev. D, 64, 043502

\bibitem[\protect\citeauthoryear{Tulin \& Yu}{2018}]{TulinYu2018}
Tulin S., Yu H.-B., 2018, Phys. Rep., 730, 1

\bibitem[\protect\citeauthoryear{Tulin et al.}{2013}]{Tulin2013}
Tulin S., Yu H.-B., Zurek K.~M., 2013, Phys. Rev. D, 87, 115007

\bibitem[\protect\citeauthoryear{Vogelsberger et al.}{2019}]{Vogelsberger2019}
Vogelsberger M., Zavala J., Schutz K., Slatyer T.~R., 2019, MNRAS, 484, 5437

\bibitem[\protect\citeauthoryear{Winther et al.}{2015}]{Winther2015}
Winther H.~A., et al., 2015, MNRAS, 454, 4208

\end{thebibliography}

\label{lastpage}
\end{document}
